% Options for packages loaded elsewhere
% Options for packages loaded elsewhere
\PassOptionsToPackage{unicode}{hyperref}
\PassOptionsToPackage{hyphens}{url}
\PassOptionsToPackage{dvipsnames,svgnames,x11names}{xcolor}
%
\documentclass[
  letterpaper,
  DIV=11,
  numbers=noendperiod]{scrartcl}
\usepackage{xcolor}
\usepackage{amsmath,amssymb}
\setcounter{secnumdepth}{-\maxdimen} % remove section numbering
\usepackage{iftex}
\ifPDFTeX
  \usepackage[T1]{fontenc}
  \usepackage[utf8]{inputenc}
  \usepackage{textcomp} % provide euro and other symbols
\else % if luatex or xetex
  \usepackage{unicode-math} % this also loads fontspec
  \defaultfontfeatures{Scale=MatchLowercase}
  \defaultfontfeatures[\rmfamily]{Ligatures=TeX,Scale=1}
\fi
\usepackage{lmodern}
\ifPDFTeX\else
  % xetex/luatex font selection
\fi
% Use upquote if available, for straight quotes in verbatim environments
\IfFileExists{upquote.sty}{\usepackage{upquote}}{}
\IfFileExists{microtype.sty}{% use microtype if available
  \usepackage[]{microtype}
  \UseMicrotypeSet[protrusion]{basicmath} % disable protrusion for tt fonts
}{}
\makeatletter
\@ifundefined{KOMAClassName}{% if non-KOMA class
  \IfFileExists{parskip.sty}{%
    \usepackage{parskip}
  }{% else
    \setlength{\parindent}{0pt}
    \setlength{\parskip}{6pt plus 2pt minus 1pt}}
}{% if KOMA class
  \KOMAoptions{parskip=half}}
\makeatother
% Make \paragraph and \subparagraph free-standing
\makeatletter
\ifx\paragraph\undefined\else
  \let\oldparagraph\paragraph
  \renewcommand{\paragraph}{
    \@ifstar
      \xxxParagraphStar
      \xxxParagraphNoStar
  }
  \newcommand{\xxxParagraphStar}[1]{\oldparagraph*{#1}\mbox{}}
  \newcommand{\xxxParagraphNoStar}[1]{\oldparagraph{#1}\mbox{}}
\fi
\ifx\subparagraph\undefined\else
  \let\oldsubparagraph\subparagraph
  \renewcommand{\subparagraph}{
    \@ifstar
      \xxxSubParagraphStar
      \xxxSubParagraphNoStar
  }
  \newcommand{\xxxSubParagraphStar}[1]{\oldsubparagraph*{#1}\mbox{}}
  \newcommand{\xxxSubParagraphNoStar}[1]{\oldsubparagraph{#1}\mbox{}}
\fi
\makeatother


\usepackage{longtable,booktabs,array}
\newcounter{none} % for unnumbered tables
\usepackage{calc} % for calculating minipage widths
% Correct order of tables after \paragraph or \subparagraph
\usepackage{etoolbox}
\makeatletter
\patchcmd\longtable{\par}{\if@noskipsec\mbox{}\fi\par}{}{}
\makeatother
% Allow footnotes in longtable head/foot
\IfFileExists{footnotehyper.sty}{\usepackage{footnotehyper}}{\usepackage{footnote}}
\makesavenoteenv{longtable}
\usepackage{graphicx}
\makeatletter
\newsavebox\pandoc@box
\newcommand*\pandocbounded[1]{% scales image to fit in text height/width
  \sbox\pandoc@box{#1}%
  \Gscale@div\@tempa{\textheight}{\dimexpr\ht\pandoc@box+\dp\pandoc@box\relax}%
  \Gscale@div\@tempb{\linewidth}{\wd\pandoc@box}%
  \ifdim\@tempb\p@<\@tempa\p@\let\@tempa\@tempb\fi% select the smaller of both
  \ifdim\@tempa\p@<\p@\scalebox{\@tempa}{\usebox\pandoc@box}%
  \else\usebox{\pandoc@box}%
  \fi%
}
% Set default figure placement to htbp
\def\fps@figure{htbp}
\makeatother


% definitions for citeproc citations
\NewDocumentCommand\citeproctext{}{}
\NewDocumentCommand\citeproc{mm}{%
  \begingroup\def\citeproctext{#2}\cite{#1}\endgroup}
\makeatletter
 % allow citations to break across lines
 \let\@cite@ofmt\@firstofone
 % avoid brackets around text for \cite:
 \def\@biblabel#1{}
 \def\@cite#1#2{{#1\if@tempswa , #2\fi}}
\makeatother
\newlength{\cslhangindent}
\setlength{\cslhangindent}{1.5em}
\newlength{\csllabelwidth}
\setlength{\csllabelwidth}{3em}
\newenvironment{CSLReferences}[2] % #1 hanging-indent, #2 entry-spacing
 {\begin{list}{}{%
  \setlength{\itemindent}{0pt}
  \setlength{\leftmargin}{0pt}
  \setlength{\parsep}{0pt}
  % turn on hanging indent if param 1 is 1
  \ifodd #1
   \setlength{\leftmargin}{\cslhangindent}
   \setlength{\itemindent}{-1\cslhangindent}
  \fi
  % set entry spacing
  \setlength{\itemsep}{#2\baselineskip}}}
 {\end{list}}
\usepackage{calc}
\newcommand{\CSLBlock}[1]{\hfill\break\parbox[t]{\linewidth}{\strut\ignorespaces#1\strut}}
\newcommand{\CSLLeftMargin}[1]{\parbox[t]{\csllabelwidth}{\strut#1\strut}}
\newcommand{\CSLRightInline}[1]{\parbox[t]{\linewidth - \csllabelwidth}{\strut#1\strut}}
\newcommand{\CSLIndent}[1]{\hspace{\cslhangindent}#1}



\setlength{\emergencystretch}{3em} % prevent overfull lines

\providecommand{\tightlist}{%
  \setlength{\itemsep}{0pt}\setlength{\parskip}{0pt}}



 


\newcommand{\doctitle}{Wiener, Kybernetik: annotations}
\newcommand{\docdate}{2026-05-01}
\newcommand{\doclogoleft}{assets/fadenlogo-fwb.png}
\newcommand{\doclogoright}{assets/dcl-head-niemand.png}
% ============================================================
% loani.tex  —  include-in-header snippet, LuaLaTeX
% IPA handled by ipa.lua Lua filter at AST level.
% No ucharclasses, no babel, no runtime font switching.
% ============================================================

\usepackage{iftex}
\RequireLuaTeX

\usepackage{fontspec}

\setmainfont{Georgia}[
  Ligatures = TeX
]
% \usepackage[hebrew]{babel}
% Hebrew font — switched via ucharclasses (XeTeX) or
% directly referenced where needed
% \newfontfamily\hebrewfont{Alef}[
%   Script = Hebrew
% ]
% \newfontfamily\babelfont{Alef}[
%   Script = Hebrew
% ]
% \usepackage[hebrew, english]{babel}
% \babelfont[hebrew]{rm}{Alef}
% \usepackage[english, bidi=basic]{babel}
% \usepackage[english, hebrew, bidi=basic]{babel}
\usepackage[english, bidi=basic]{babel}

% \babelprovide[onchar=fonts ids, main=false]{hebrew} % all RTL, IPA wrongs
% \babelprovide[onchar=ids fonts, main=false]{english}
% \babelprovide[onchar=fonts ids, main=true]{english}
\babelprovide[onchar=fonts ids]{hebrew}

\babelfont[hebrew]{rm}{Alef}
% IPA font — called explicitly by ipa.lua filter output
\newfontfamily\ipafont{Source Sans 3}

% ------------------------------------------------------------
% Hebrew bidi — luatexja-style or direct LuaTeX primitive.
% Since babel is not loaded, we use the LuaTeX direction
% primitive directly: TRT = right-to-left, TLT = left-to-right
% Hebrew Unicode chars get RTL from the UBA via luaotfload.
% ------------------------------------------------------------
\directlua{
  luatexbase.add_to_callback("pre_linebreak_filter",
    function(head) return head end,
    "dummy_bidi")
}

\usepackage[hyperfootnotes=false]{hyperref}
% ============================================================
% header.tex  —  include-in-header snippet
% Replaces $var$ Pandoc syntax with \ifdefined\doc... guards.
% Vars injected via vars.lua Lua filter from YAML front matter:
%   \doctitle      \docsubtitle   \docauthor   \docdate
%   \docmeta       \docfoot       \doclogoleft \doclogoright
%
% No fancyhdr. Header fires via \AtBeginDocument,
% footer via \AtEndDocument — bidi-safe.
% ============================================================

\usepackage{graphicx}
\usepackage{tabularx}
\usepackage{xcolor}
\usepackage{array}
\usepackage{tikz}

% ---------- header font (adjust to taste) ----------
\newfontfamily\headerfont{Source Sans 3}

% ---------- footer colors ----------
\definecolor{footerbg}{HTML}{000000}
\definecolor{footerfg}{HTML}{FFFFFF}

% ---------- column types ----------
\newcolumntype{L}{>{\raggedright\arraybackslash}m{10mm}}
\newcolumntype{C}{>{\centering\arraybackslash}X}
\newcolumntype{R}{>{\raggedleft\arraybackslash}m{45mm}}

% vertical separator rule
\newcommand{\vsep}{\color{black}\vrule width 1pt}

% ============================================================
% CUSTOM HEADER
% ============================================================
\newcommand{\CustomHeader}{%
% \LTR{%
\noindent
\begin{tabularx}{\textwidth}{L !{\vsep} C !{\vsep} R}

% LEFT LOGO
\ifdefined\doclogoleft
  \includegraphics[width=10mm,height=10mm,keepaspectratio]{\doclogoleft}%
\fi
&
% CENTER: title + subtitle stacked in a minipage so \\ is a
% line break within the cell, not a new table row
\begin{minipage}[c]{\linewidth}
\centering
\ifdefined\doctitle
  {\Large\bfseries\headerfont\doctitle}%
\fi
\ifdefined\docsubtitle
  \\[2pt]{\normalsize\headerfont\docsubtitle}%
\fi
\vspace{5pt}

\end{minipage}
% &
% % CENTER TITLE BLOCK
% \ifdefined\doctitle
%   {\Large\bfseries\doctitle}%
% \fi
% \ifdefined\docsubtitle
%   \\\normalsize\docsubtitle
% \fi
&
% RIGHT LOGO
\ifdefined\doclogoright
  \includegraphics[height=10mm,keepaspectratio]{\doclogoright}%
\fi

\end{tabularx}

\hrule
% \vspace{6pt}

% ---------- META AREA ----------
{\small
\ifdefined\docauthor
  \textbf{Author:} \docauthor\\
\fi
\ifdefined\docdate
  \textbf{Date:} \docdate
\fi
\ifdefined\docmeta
  \\\docmeta\\
  snc: 16132.  
  {\ipafont ɪ ʃ ɔ ː ʁ}  
  {\ipafont wer es [pɛʟɛχɪʟ] ausspricht, und nicht [pɛʀɛχɪʟ]}
\fi
}
\vspace{6pt}

\hrule

% \vspace{12pt}
% }% end \LTR
}% end \CustomHeader

% ============================================================
% CUSTOM FOOTER  (overlay, anchored to page bottom)
% ============================================================
\newcommand{\CustomFooter}{%
\begin{tikzpicture}[remember picture, overlay]
\node[
  anchor=south,
  draw,
  fill=footerbg,
  text=footerfg,
  inner sep=10pt,
  rounded corners=2pt,
  text width=\dimexpr\textwidth-20pt\relax
]
at ([yshift=5mm]current page.south)
{%
  \begin{center}
  {\small
  \ifdefined\docfoot
    % {\docfoot}%
    {\headerfont\docfoot}%
  \fi
  }
  \end{center}
};
\end{tikzpicture}%
}% end \CustomFooter

% ============================================================
% FIRE header at doc start, footer at doc end
% ============================================================
% \AtBeginDocument{\CustomHeader}
\AtEndDocument{\CustomFooter}
\renewcommand{\maketitle}{}
\usepackage{booktabs}
\usepackage{longtable}
\usepackage{array}
\usepackage{multirow}
\usepackage{wrapfig}
\usepackage{float}
\usepackage{colortbl}
\usepackage{pdflscape}
\usepackage{tabu}
\usepackage{threeparttable}
\usepackage{threeparttablex}
\usepackage[normalem]{ulem}
\usepackage{makecell}
\usepackage{xcolor}
\KOMAoption{captions}{tableheading}
\makeatletter
\@ifpackageloaded{caption}{}{\usepackage{caption}}
\AtBeginDocument{%
\ifdefined\contentsname
  \renewcommand*\contentsname{Table of contents}
\else
  \newcommand\contentsname{Table of contents}
\fi
\ifdefined\listfigurename
  \renewcommand*\listfigurename{List of Figures}
\else
  \newcommand\listfigurename{List of Figures}
\fi
\ifdefined\listtablename
  \renewcommand*\listtablename{List of Tables}
\else
  \newcommand\listtablename{List of Tables}
\fi
\ifdefined\figurename
  \renewcommand*\figurename{Figure}
\else
  \newcommand\figurename{Figure}
\fi
\ifdefined\tablename
  \renewcommand*\tablename{Table}
\else
  \newcommand\tablename{Table}
\fi
}
\@ifpackageloaded{float}{}{\usepackage{float}}
\floatstyle{ruled}
\@ifundefined{c@chapter}{\newfloat{codelisting}{h}{lop}}{\newfloat{codelisting}{h}{lop}[chapter]}
\floatname{codelisting}{Listing}
\newcommand*\listoflistings{\listof{codelisting}{List of Listings}}
\makeatother
\makeatletter
\makeatother
\makeatletter
\@ifpackageloaded{caption}{}{\usepackage{caption}}
\@ifpackageloaded{subcaption}{}{\usepackage{subcaption}}
\makeatother
\usepackage{bookmark}
\IfFileExists{xurl.sty}{\usepackage{xurl}}{} % add URL line breaks if available
\urlstyle{same}
\hypersetup{
  colorlinks=true,
  linkcolor={blue},
  filecolor={Maroon},
  citecolor={Blue},
  urlcolor={Blue},
  pdfcreator={LaTeX via pandoc}}


\title{Wiener, Kybernetik: annotations}
\author{}
\date{2026-05-01}
\begin{document}
\maketitle

\CustomHeader


\subsection{notes on papers}\label{notes-on-papers}

\subsubsection{snc}\label{snc}

\begin{itemize}
\tightlist
\item
  t1
\end{itemize}

\paragraph{annotations:
Foerster\_Ethik+und+Kybernetik+zweiter+Ordnung\_m
Kopie.pdf}\label{annotations-foerster_ethikundkybernetikzweiterordnung_m-kopie.pdf}

paper: Foerster (1993)

page

note

comment

1

Heinz von FoersterKybernEthik

\#cite\_ Foerster\_Ethik+und+Kybernetik+zweiter+Ordnung\_m Kopie.pdf\_\_

1

Heinz von FoersterKybernEthik

\#cite\_ Foerster\_Ethik+und+Kybernetik+zweiter+Ordnung\_m Kopie.pdf\_\_

4

der Kybernetiker, Gordon Pask: ``Kybernetik ist die Wissenschaft von
vertretbaren Metaphern

5

Im allgemeinen Fall des zirkulären Schlusses bedeutet A impliziert B; B
impliziert C; und - zum allgemeinen Entsetzen - C impliziert A! Oder, im
reflexiven Fall: A impliziert B; und - Oh, Grauen! - B impliziert A

5

NA

marginnote4app://note/CF9607AC-C9E6-424F-8676-CBDC3374F847

7

``Der erste Gedanke bei der Aufstellung eines ethischen Gesetzes von der
Form `Du sollst\ldots{}' ist: Und was dann, wenn ich es nicht tue

8

Oder denken wir an Goldbachs ``Vermutung'', die sich derartig einfach
anhört, daß man glaubt, ein Beweis liege fast auf der Hand: ``Jede
gerade Zahl ist die Summe zweier Primzahlen.'' Zum Beispiel: 12 ist die
Summe der zwei Primzahlen 5

11

das Problem liegt im Verstehen des Verstehens; das Problem besteht
darin, Entscheidungen über prinzipiell unentscheidbare Fragen zu treffen

12

Ethik antwortete: ``Sag ihnen, sie sollten immer so handeln, die Anzahl
der Möglichkeiten zu vermehren; ja, die Anzahl der Möglichkeiten zu
vermehren

\paragraph{annotations:
Wiener\_Einführung.pdf}\label{annotations-wiener_einfuhrung.pdf}

paper: Wiener (1992)

page

note

comment

5

EINFÜHRUNG

\#init\_Wiener, Kybernetik\_\_(..)\#cite\_wiener\_kybernetik\_1992\_\_

5

EINFÜHRUNG

\#init\_Wiener, Kybernetik\_\_(..)\#cite\_wiener\_kybernetik\_1992\_\_

6

Viele Jahre hatten Dr.~Rosenblueth und ich die Überzeugung geteilt, daß
die für das Gedeihen der Wissenschaft fruchtbarsten Gebiete jene waren,
die als Niemandsland zwischen den verschiedenen bestehenden Disziplinen
vernachlässigt wurden. Seit Leibniz hat es vielleicht keinen Menschen
mehr gegeben, der die volle Übersicht über die gesamte geistige
Tätigkeit seiner Zeit gehabt hat. (..)(..)Seit jener Zeit ist die
Wissenschaft in zunehmendem Maß die Aufgabe von Spezialisten geworden;
auf Gebieten, die die Tendenz zeigen, ständig näher zusammenzuwachsen

7

Wir haben jahrelang von einem Institut mit unabhängigen Wissenschaftlern
geträumt, die gemeinsam in diesen Hinterwäldern der Wissenschaft
arbeiten würden; nicht als Untergeordnete irgendeines hohen
Exekutivbeamten, sondern vereint durch den Wunsch --- ja durch die
geistige Notwendigkeit ---, das Teilgebiet als Ganzes zu verstehen und
einander zu diesem Verstehen zu verhelfen

8

Bei Kriegsbeginn richteten das deutsche Luftwaffenpotential und die
defensive Lage Englands die Aufmerksamkeit vieler Wissenschaftler auf
die Entwicklung der Flugabwehrartillerie. Schon vor dem Krieg war es
klargeworden, daß die Geschwindigkeit des Flugzeugs alle klassischen
Methoden der Feuerleitung überwunden hatte und daß es nötig war, alle
notwendigen Rechnungen in die Regelungsapparatur selbst einzubauen.
Diese waren sehr schwierig geartet durch die Tatsache, daß --- nicht zu
vergleichen mit allen vorher betrachteten Zielen --- ein Flugzeug eine
Geschwindigkeit hat, die ein sehr ansehnlicher Bruchteil der
Geschwindigkeit des Geschosses ist, das zum Beschuß verwendet wird.
Demgemäß ist es außerordentlich wichtig, das Geschoß nicht auf das Ziel
abzuschießen, sondern so, daß Geschoß und Ziel im Raum zu einem späteren
Zeitpunkt zusammentreffen. (..)(..)Wir mußten deshalb eine Methode
Enden, die zukünftige Position des Flugzeuges vorherzusagen

9

ein außerordentlich wichtiger Faktor im willensgesteuerten Handeln das
ist, was die Regelungstechniker mit Ruckkopplung bezeichnen

9

Hier geniigt die Feststellung, daB bei einer von einem Muster gelenkten
Bewegung die Abweichung der wirklich durchgeführten Bewegung von diesem
Muster als neue Eingabe benutzt wird, um den geregelten Teil zu
veranlassen, die Bewegung dem Muster naherzubringen.

9

Nehmen wir nun an, daB ich einen Bleistift aufhebe, so muB ich, um
dieses zu tun, gewisse Muskeln bewegen. Jedoch jeder von uns,
ausgenommen wenige Anatomieexperten, weiß nicht, welches diese Muskeln
sind, und sogar unter den Anatomen gibt es wenige, wenn es ilberhaupt
welche gibt, die diese Handlung durch eine bewußte Willenssteuerung der
Kontraktion jedes betreffenden Muskels ausführen können. Was wir
hingegen wollen, ist, den Bleistift aufzuheben. Haben wir dies einmal
beschlossen, so schreitet unsere Bewegung derart fort, daß wir grob
sagen können, daß der Grad, zu welchem der Bleistift noch nicht
aufgehoben ist, mit jedem Bewegungsstadium vermindert wird. Dieser Teil
der Aktion ist nicht voll im Bewußtsein.

10

Die Nachricht ist eine zeitlich diskret oder stetig verteilte Folge
meßbarer Ereignisse --- genau das, was von den Statistikern ein
Zufallsprozeß genannt wird. Die Vorhersage der Zukunft einer Nachricht
geschieht durch irgendeine Operation auf ihre Vergangenheit,
gleichgültig, ob dieser Operator durch ein mathematisches Rechenschema
oder durch einen mechanischen oder elektrischen Apparat verwirklicht
wird.

11

Dieses sich gegenseitig beeinflussende Paar von Fehlertypen schien etwas
mit dem kontrastierenden Problem der Ortsmessung und der
Bewegungsmessung zu tun zu haben, das in der Quantenmechanik von
Heisenberg als Ungenauigkeitsrelation zu finden ist.

1927

11

Dieses sich gegenseitig beeinflussende Paar von Fehlertypen schien etwas
mit dem kontrastierenden Problem der Ortsmessung und der
Bewegungsmessung zu tun zu haben, das in der Quantenmechanik von
Heisenberg als Ungenauigkeitsrelation zu finden ist.

1927

12

NA

\#bookmarks Page 12

12

Governor« für Fliehkraftregler ist von einer lateinischen Verfälschung
von κυβερνήτης abgeleitet

12

Wir haben beschlossen, das ganze Gebiet der Regelung und
Nachrichtentheorie, ob in der Maschine oder im Tier, mit dem Namen
»Kybernetik« zu benennen, den wir aus dem griechischen
»κvßερvήΤη\textasciitilde« oder »Steuermann« bildeten

12

Informationsgehaltes berührt in natürlicher Weise einen klassischen
Begriff in der statistischen Mechanik: den der Entropie. Gerade wie der
Informationsgehalt eines Systems ein Maß des Grades der Ordnung ist, ist
die Entropie eines Systems ein MaB des Grades der Unordnung; und das
eine ist einfach das Negative des anderen.

13

Gruppe viel dazu beigetragen hat, die Aufmerksamkeit der mathematisch
Interessierten auf die Möglichkeiten der biologischen Wissenschaften
hinzulenken,

13

Es ist deshalb nicht im mindesten überraschend, daB der gleiche
intellektuelle Impuls, der zur Entwicklung der mathematischen Logik
geführt hat, gleichzeitig zur idealen oder tatsächlichen Mechanisierung
der Prozesse des Denkens geführt hat.

13

Vereinigung der Nervenfasern durch Synapsen zu Systemen mit gegebenen
Gesamteigenschaften

14

NA

\#bookmarks Page 14

14

Der Alles-oder-nichts-Charakter der Neuronenentladung ist völlig analog
zur Auswahl einer binären Ziffer; und schon mehr als einer von uns hatte
das binäre Zahlensystem als beste Basis des Rechnens in der Maschine
erkannt. Die Synapse ist nichts als ein Mechanismus, der bestimmt, ob
eine gewisse Kombination von Ausgängen von anderen Elementen ein
ausreichender Anreiz für das Entladen des nächsten Elementes ist oder
nicht und muß ein genaues Analogon in der Rechenmaschine haben

14

Beispiele von modernen Vakuumröhren zeigte und ihm erklärte, daß diese
ideale Mittel wären, um apparative Aquivalente zu seinen nervlichen
Kreisen und Systemen darzustellen

14

ultraschnelle Rechenmaschine, so wie sie abhängig war von
aufeinanderfolgenden Schaltern, beinahe ein ideales Modell der sich aus
dem Nervensystem ergebenden Probleme darstellen mußte

15

Die Physiologen gaben eine gemeinsame Darstellung von
Kybernetikproblemen von ihrem Gesichtspunkt aus; ähnlich stellten die
Rechenmaschinenbauer ihre Methoden und Ziele dar. Am Ende des Treffens
war es allen klar, daß es eine beträchtliche gemeinsame Denkbasis aller
Bearbeiter der verschiedenen Gebiete gab, daB man in jeder Gruppe schon
Begriffe gebrauchen konnte, die durch andere schon besser entwickelt
waren, und daß ein Versuch gemacht werden sollte, ein allgemeines
Vokabular zustande zu bringen.

15

So waren auf diese oder jene Weise bei Kriegsende die Gedanken der
Vorhersagetheorie und der statistischen Nachrichtentheorie bereits einem
großen Teil der Statistiker und Nachrichteningenieure der Vereinigten
Staaten und Großbritanniens bekannt.

16

Forschung schlug zwei Richtungen ein: das Studium der Phänomene der
Leitfähigkeit und Latenz in gleichförmig leitenden zwei- oder
dreidimensionalen Medien und die statistische Untersuchung der
Leitungseigenschaften von zufälligen Netzen aus leitenden Fasern.

16

Vieles von der früheren Psychologie hat sich als nichts anderes als die
Physiologie der speziellen Sinnesorgane herausgestellt, und das gesamte
Gewicht des Gedankengutes, das die Kybernetik in die Psychologie
hineinträgt, betrifft die Physiologie und Anatomie

16

daß der Herzmuskel ein reizbares Gewebe darstellte, das ebenso für die
Erforschung von Leitungsmechanismen brauchbar war wie das Nervengewebe

16

Der Kern unserer Treffen war die Gruppe gewesen, die in Princeton 1944
versammelt war, aber die Doktoren McCulloch und FremontSmith hatten
richtig die psychologischen und soziologischen Beziehungen des Themas
erkannt und der Gruppe eine Anzahl von führenden Psychologen, Soziologen
und Anthropologen hinzugeladen

17

Von Anfang an haben wir vorausgesehen, daß das Problem des Erkennens von
»Gestalter« oder der wahrnehmbaren Formation der Allgemeinbegriffe zu
dieser Art gehδren würde. Was ist der Mechanismus, durch den wir ein
Quadrat als ein Quadrat erkennen, ohne Rllcksicht auf seine Lage, seine
Größe und seine Orientierung

17

erstes Treffen, das im Frühling 1946 abgehalten wurde, war größtenteils
didaktischen Vorträgen für diejenigen von uns gewidmet, die beim
Princeton-Treffen dabeigewesen waren, und einer allgemeinen Festlegung
der Bedeutung des Gebietes für alle Anwesenden

17

Was die Soziologie und Anthropologie betrifft, ist es offenkundig, daB
die Information und Übertragung als Mechanismus der fortschreitenden
Organisation vom Einzelwesen zur Gemeinschaft von Bedeutung ist

17

ein nervliches Problem direkt mit dem Begriff der Rückkopplung zu
behandeln

18

Dieses näherungsweise logarithmische Verhalten des Synapsenmechanismus
hängt sicher damit zusammen, daß das Weber-Fechnersche Gesetz der
Reizintensität logarithmisch ist, obgleich dieses Gesetz nur eine erste
Näherung darstellt.

19

McCulloch war vor das Problem gestellt worden, einen Apparat zu
konstruieren, der den Blinden in die Lage versetzen sollte, die
gedruckte Seite mit Hilfe des Ohres zu lesen

20

Es ist bestimmt so, daB das soziale System eine Organisation ähnlich dem
Einzelwesen ist

21

Es gibt zwei andere Gebiete, wo ich letztlich hoffe, einiges mit Hilfe
kybernetischer Gedanken tatsächlich zustande zu bringen, bei denen diese
Hoffnung jedoch auf weitere Entwicklungen warten muß. Eines davon
betrifft Prothesen für verlorene oder verkümmerte Glieder

22

Es war mir schon lange klar gewesen, daß die modernen, ultraschnellen
Rechenmaschinen im Prinzip ein ideales zentrales Nervensystem für eine
automatische Regelungsanlage waren und daB ihre Eingaben und Ausgänge
nicht die Form von Zahlen oder Diagrammen haben müssen, sondern sehr gut
Äußerungen von künstlichen Sinnesorganen sein können, wie z. B. von
photoelektrischen Zellen oder Thermometern bzw. die Wirkungen von
Motoren oder Magnetspulen. Mit der Hilfe von Spannungsmessern oder
ähnlichen Hilfsmitteln, die die Wirkungen dieser Motororgane lesen und
berichten und zu dem zentralen Steuerungssystem als einem künstlichen
kinästhetischen Sinn zurückkoppeln, sind wir bereits in der Lage,
künstliche Maschinen von fast jedem Grad sorgfältiger Arbeitsleistung zu
konstruieren

23

Diejenigen von uns, die zu der neuen Wissenschaft Kybernetik beigetragen
haben, sind in einer moralischen Lage, die, um es gelinde auszudrücken,
nicht sehr bequem ist. Wir haben zu der Einführung einer neuen
Wissenschaft beigesteuert, die, wie ich gesagt habe, technische
Entwicklungen mit großen Möglichkeiten für Gut oder Böse umschließt. Wir
können sie nur in die Welt weitergeben, die um uns existiert, und dies
ist die Welt von Belsen und Hiroshima

23

Natürlich, gerade wie der gelernte Zimmermann, der gelernte Mechaniker,
der gelernte Schneider in gewissem Grade die erste industrielle
Revolution überlebt haben, können der erfahrene Wissenschaftler und der
erfahrene Verwaltungsbeamte die zweite überleben. Wenn man sich jedoch
die zweite Revolution abgeschlossen denkt, hat das durchschnittliche
menschliche Wesen mit mittelmäßigen oder noch geringeren Kenntnissen
nichts zu verkaufen, was für irgend jemanden das Geld wert wäre

23

Die moderne industrielle Revolution ist in ähnlicher Weise dazu
bestimmt, das menschliche Gehirn zu entwerten, wenigstens in seinen
einfacheren und mehr routinemäßigen Entscheidungen

23

Es kann sehr wohl fΥr die Menschheit gut sein, Maschinen zu besitzen,
die sie von der Notwendigkeit niedriger und unangenehmer Aufgaben
befreien, oder es kann auch nicht gut sein. Ich weiß es nicht. Es kann
nicht gut sein, diese neuen Kräfteverhältnisse in den Begriffen des
Marktes abzuschätzen, des Geldes, das sie verdienen; und es sind genau
die Begriffe des freien Marktes, des »fifth freedom«, die die
Losungsworte des Teiles der amerikanischen Meinung wurden, der durch die
National Association of Manufacturers und die Saturday Evening Post
repräsentiert wird

24

Das Beste, was wir tun können, ist, zu versuchen, daB eine breite
Öffentlichkeit die Richtung und die Lage der gegenwärtigen Arbeit
versteht, und unsere persönlichen Anstrengungen auf die Gebiete zu
beschränken, die, wie z. B. Physiologie und Psychologie, am weitesten
von Krieg und Unterdrückung entfernt sind.

\protect\phantomsection\label{refs}
\begin{CSLReferences}{1}{1}
\bibitem[\citeproctext]{ref-foerster_kybernethik_1993}
Foerster, Heinz. 1993. \emph{{KybernEthik}}. Internationaler
{Merve}-{Diskurs} 180. Merve.

\bibitem[\citeproctext]{ref-wiener_kybernetik_1992}
Wiener, Norbert. 1992. \emph{Kybernetik: {Regelung} Und
{Nachrichtenübertragung} Im {Lebewesen} Und in Der {Maschine} ({MIT},
1948)}. 2nd ed. Econ Classics. ECON-Verl.

\end{CSLReferences}




\end{document}
